\section*{Chapter \ref{ch:g12:metamorphic-rocks} --- Metamorphic Rocks}

\begin{solution}{exo:g12:metamorphic-rocks:1}
(a) \omterm{def:g12:metamorphic-rocks:polymorph}{polymorphs}; (b) an \omterm{def:g12:metamorphic-rocks:index-mineral}{isograd}; (c) an \omterm{def:g12:metamorphic-rocks:contact-regional}{aureole}; (d) a \omterm{def:g12:metamorphic-rocks:foliation}{foliation};
(e) the \omterm{def:g12:metamorphic-rocks:lithostatic}{lithostatic pressure}.
\end{solution}

\begin{solution}{exo:g12:metamorphic-rocks:2}
(1) marble; (2) quartzite; (3) slate; (4) gneiss; (5) hornfels.
\end{solution}

\begin{solution}{exo:g12:metamorphic-rocks:3}
$P = \rho g z$: at \qty{10}{km}, $2830 \times 9.81 \times \num{10000} \approx
\qty{2.78e8}{Pa} = \qty{0.278}{GPa}$; at \qty{20}{km}, twice as much,
\qty{0.555}{GPa}.
\end{solution}

\begin{solution}{exo:g12:metamorphic-rocks:4}
(1) Kyanite (the kyanite--andalusite line is near \qty{0.25}{GPa} at
\qty{400}{\celsius}). (2) Andalusite. (3) Sillimanite.
\end{solution}

\begin{solution}{exo:g12:metamorphic-rocks:5}
(1) Slate, phyllite, schist, gneiss. (2) In the schist. (3) The gneiss.
\end{solution}

\begin{solution}{exo:g12:metamorphic-rocks:6}
(1) $2830 \times 9.81 \times \num{35000} \approx \qty{9.72e8}{Pa}$, that is
$\num{9.72e8}/101325 \approx \num{9600}$ atmospheres. (2) $1025 \times 9.81
\times 4000 \approx \qty{4.02e7}{Pa} = \qty{40.2}{MPa}$. (3) $972/40.2
\approx 24$ times.
\end{solution}

\begin{solution}{exo:g12:metamorphic-rocks:7}
(1) $z = \num{0.375e9}/(2830 \times 9.81) \approx \qty{13.5}{km}$.
(2) $(504 - 10)/13.5 \approx \qty{36.6}{\celsius}$ per kilometre. (3) No:
the rock would have to sit exactly at one point of the diagram; such rocks
are rare.
\end{solution}

\begin{solution}{exo:g12:metamorphic-rocks:8}
(1) $23/3 \approx \qty{8}{mm}$. (2) About 8 million years. (3) The rock
stayed at depth, at a high temperature, for millions of years.
\end{solution}

\begin{solution}{exo:g12:metamorphic-rocks:9}
(1) Four: the garnet, staurolite, kyanite and sillimanite \omterm{def:g12:metamorphic-rocks:index-mineral}{isograds}, in that
order. (2) Westward. (3) At the western end, in the sillimanite zone. Its
rocks crossed the kyanite--sillimanite line: buried about \qty{20}{km}
(\qty{0.56}{GPa}), they were hotter than about \qty{590}{\celsius}.
\end{solution}

\begin{solution}{exo:g12:metamorphic-rocks:10}
(1) Mud becomes hornfels, limestone becomes marble, sandstone becomes
quartzite. (2) They were heated without being squeezed: no directed pressure
aligned their minerals. (3) The rocks around it are already hot, so the
magma heats them little more.
\end{solution}

\begin{solution}{exo:g12:metamorphic-rocks:11}
(1) Across the push: their flat faces face east and west, so the flakes
stand in north--south planes. (2) Vertical, running north--south.
(3) Calcite and quartz grains are neither flat nor long, so even when they
recrystallise under a squeeze there is nothing to line up.
\end{solution}

\begin{solution}{exo:g12:metamorphic-rocks:12}
(1) $z = (T - 10)/30$, so $P = 0.0278\,(T - 10)/30 \approx
0.000927\,(T - 10)$. (2) $0.000927\,(T - 10) = 0.577 + 0.00209\,(T - 600)$
gives $0.001163\,T \approx 0.668$, so $T \approx \qty{574}{\celsius}$, $P
\approx 0.000927 \times 564 \approx \qty{0.523}{GPa}$ and $z = (574 - 10)/30
\approx \qty{18.8}{km}$. (3) Kyanite above that depth, sillimanite below it.
\end{solution}

\begin{solution}{exo:g12:metamorphic-rocks:13}
(1) $10 + 10 \times 35 = \qty{360}{\celsius}$. (2) $10 + 30 \times 35 =
\qty{1060}{\celsius}$. (3) The pressure at \qty{35}{km} is about
\qty{0.97}{GPa}; at \qty{360}{\celsius} that point lies deep in the kyanite
field. The second point lies far to the right of the diagram, in the
sillimanite field, at temperatures where crustal rocks begin to melt.
(4) They need a high pressure at a low temperature, a combination that only
the cold crust carried down in a \omterm{def:g11:subduction-zones:subduction}{subduction zone} provides.
\end{solution}

\begin{solution}{exo:g12:metamorphic-rocks:14}
(1) Kyanite $1000/3.61 \approx \qty{277}{cm^3}$; sillimanite $1000/3.24
\approx \qty{309}{cm^3}$; andalusite $1000/3.15 \approx \qty{317}{cm^3}$.
(2) $(317 - 277)/317 \approx 0.127$, about \qty{13}{\percent}. (3) At high
pressure the most compact arrangement of the atoms is favoured; kyanite
packs the same atoms into the least room.
\end{solution}

\begin{solution}{exo:g12:metamorphic-rocks:15}
(1) \qty{0.555}{GPa} at \qty{20}{km}, \qty{0.694}{GPa} at \qty{25}{km}.
(2) $T = 600 + (P - 0.577)/0.00209$: about \qty{590}{\celsius} at
\qty{0.555}{GPa} and \qty{656}{\celsius} at \qty{0.694}{GPa}. At
\qty{700}{\celsius} the rocks were hotter than both, in the sillimanite
field: consistent. (3) $(700 - 10)/20 \approx \qty{34.5}{\celsius}$ and
$(700 - 10)/25 \approx \qty{27.6}{\celsius}$ per kilometre, close to the
typical gradient. (4) $\qty{20}{km}$ to $\qty{25}{km}$ in 40 million years:
0.50 to \qty{0.63}{km} per million years, that is 0.50 to \qty{0.63}{mm} a
year.
\end{solution}

\begin{solution}{pb:g12:metamorphic-rocks:1}
\textbf{1.} It is foliated, and its micas are big enough to be seen. The
micas, muscovite and biotite, make the \omterm{def:g12:metamorphic-rocks:foliation}{schistosity}.

\textbf{2.} \ce{Al2SiO5}; they are \omterm{def:g12:metamorphic-rocks:polymorph}{polymorphs}.

\textbf{3.} The volumes of a kilogram are in the ratio $3.15/3.61 \approx
0.873$: kyanite is about \qty{13}{\percent} smaller. Pressure favours
kyanite.

\textbf{4.} The kyanite zone. On its way down it crossed the biotite,
garnet, staurolite and kyanite \omterm{def:g12:metamorphic-rocks:index-mineral}{isograds}.

\textbf{5.} About 8 million years.

\textbf{6.} $P = 0.577 + 0.00209 \times 0 = \qty{0.577}{GPa}$.

\textbf{7.} \qty{0.70}{GPa} is above \qty{0.577}{GPa}, on the kyanite side of
the kyanite--sillimanite line; it is also far above the andalusite field,
which at \qty{600}{\celsius} lies below about \qty{0.23}{GPa}.

\textbf{8.} \qty{0.577}{GPa}.

\textbf{9.} $0.70 = 0.577 + 0.00209\,(T - 600)$, so $T = 600 +
0.123/0.00209 \approx \qty{659}{\celsius}$.

\textbf{10.} \qty{7.0}{kbar}; $\num{7.0e8}/101325 \approx \num{6900}$
atmospheres.

\textbf{11.} $\rho g = 2830 \times 9.81 \approx \qty{27800}{Pa}$ per metre,
\qty{27.8}{MPa} per kilometre.

\textbf{12.} $z = \num{0.70e9}/\num{27762} \approx \qty{25200}{m}$, about
\qty{25.2}{km}.

\textbf{13.} $\num{0.577e9}/\num{27762} \approx \qty{20.8}{km}$.

\textbf{14.} $\num{0.70e9}/(2700 \times 9.81) \approx \qty{26.4}{km}$, about
\qty{5}{\percent} deeper ($26.4/25.2 \approx 1.048$).

\textbf{15.} In the crust: \qty{25}{km} is less than even a normal crustal
thickness.

\textbf{16.} $(600 - 10)/25.2 \approx \qty{23.4}{\celsius}$ per kilometre.

\textbf{17.} It is much warmer than a \omterm{def:g11:subduction-zones:subduction}{subduction zone} (under
\qty{10}{\celsius} per kilometre) and a little colder than a typical crust
(\qty{30}{\celsius}), far from a volcanic area (40 to \qty{50}{\celsius}): a
rock buried deep in the thickened crust of a mountain belt, by \omterm{def:g12:metamorphic-rocks:contact-regional}{regional
metamorphism}.

\textbf{18.} $10 + 30 \times 25.2 \approx \qty{766}{\celsius}$. The
kyanite--sillimanite line, continued, lies at $0.577 + 0.00209 \times 166
\approx \qty{0.92}{GPa}$ at that temperature; \qty{0.70}{GPa} is below it,
in the sillimanite field. A typical \omterm{def:g11:lithosphere-asthenosphere:geotherm}{geotherm} would have made sillimanite.

\textbf{19.} $\qty{25.2}{km} = \qty{2.52e7}{mm}$; at \qty{0.5}{mm} a year,
about 50 million years.

\textbf{20.} As the rock rose it cooled, so its reactions slowed down, and
the water that the reverse reactions need had been driven off on the way
down.

\textbf{21.} The kyanite schist crystallised about \textbf{\qty{25}{km}
down}, at about \qty{600}{\celsius}.
\end{solution}
